\documentclass[11pt,a4paper]{article}
\usepackage[utf8]{inputenc}\usepackage[T1]{fontenc}\usepackage[italian,english]{babel}
\usepackage{amsmath,amssymb,amsthm}\usepackage[margin=2.2cm]{geometry}\usepackage{longtable,booktabs,array,calc}
\usepackage{listings,xcolor}\usepackage{hyperref}\usepackage{url}
\providecommand{\tightlist}{\setlength{\itemsep}{0pt}\setlength{\parskip}{0pt}}
\lstset{basicstyle=\ttfamily\scriptsize,breaklines=true,frame=single,captionpos=b,columns=fullflexible,keepspaces=true,inputencoding=utf8,extendedchars=true}
\title{Proof Pursuit --- Problem 4, part 4\\ \large prove, decisioni degli agenti, codice verificato, letteratura}
\author{Team Proof Pursuit (Researcher: T. Tumini; Referee: gabundos; Creative: M. Renzi; gio3r3gio)}
\date{\today}
\begin{document}\maketitle\tableofcontents
\section*{Come leggere questo rapporto}
Per ogni cella: la richiesta ufficiale, lo stato dichiarato dal team (mai ``risolto'' senza revisione umana), la consegna
con la prova, poi la traccia delle decisioni degli agenti (Researcher: famiglia e motivo; Referee: verdetto ed errore
fatale; Creative: quando e perché è intervenuto) e il codice su cui il tentativo poggia con l'esito della riesecuzione
fidata. DIMOSTRATO, CITATO, VERIFICATO SU INTERVALLO FINITO e CONGETTURA restano distinti. L'architettura del sistema
è descritta nel README del repository.
\section{Problem 4 — Disjoint congruence classes and large gcd}
\subsection{Cella 4 (5 punti) — stato: parziale}
\subparagraph{\texorpdfstring{Parte 4 (C4) --- Every \(k\) up to
12}{Parte 4 (C4) --- Every k up to 12}}\label{parte-4-c4-every-k-up-to-12}

\textbf{Punteggio:} 5 points · \textbf{Valutazione:} Judged

A longer range, and a precise account of your method at the sizes just
beyond it. Prove the statement for every \(k \le 12\). Then, for each
\(k\) from \(9\) to \(16\) in turn, report exactly what your method
leaves undecided at that \(k\): if nothing, say so and prove it; if
something, exhibit it in full. If that disagrees with any source you
consulted, say which of the two is right and why. An answer that reports
agreement with a source it did not test will be marked wrong.

\textbf{Enunciato protetto dato agli agenti (state.json).} \texttt{Prove the statement for every k \textless{}= 12; then for each k = 9..16 report exactly what the method leaves undecided (nothing: prove it; something: exhibit it in full); compare with sources only if tested.}

\textbf{Evidenza registrata in STATUS.md.} bozza parziale in inglese in submission/ (parziali.py)

\subsubsection{Consegna (prova)}
\paragraph{Problema 4 --- Parte 4 --- bozza di consegna
PARZIALE}\label{problema-4-parte-4-bozza-di-consegna-parziale}

\textbf{Stato dichiarato: PARZIALE.} Nessuna soluzione completa: qui
sotto ciò che è stabilito, la formalizzazione, la posizione rispetto
alla letteratura e ciò che resta aperto. Nulla è dichiarato dimostrato
oltre quanto scritto.

\subparagraph{1. Risultato e ambito}\label{risultato-e-ambito}

\textbf{Richiesta ufficiale.} \#\# Parte 4 (C4) --- Every \(k\) up to 12

\textbf{Punteggio:} 5 points · \textbf{Valutazione:} Judged

A longer range, and a precise account of your method at the sizes just
beyond it. Prove the statement for every \(k \le 12\). Then, for each
\(k\) from \(9\) to \(16\) in turn, report exactly what your method
leaves undecided at that \(k\): if nothing, say so and prove it; if
something, exhibit it in full. If that disagrees with any source you
consulted, say which of the two is right and why. An answer that reports
agreement with a source it did not test will be marked wrong.

\textbf{Cosa consegniamo.} Come parte 3, esteso a \(k\le12\)
(\(L_{12}=27720\), 96 divisori); non eseguito.

\subparagraph{2. Dimostrazione}\label{dimostrazione}

Stessa riduzione e stesso schema di ricerca; il rapporto sui casi
\(9\le k\le16\) richiede di elencare i sopravvissuti alla potatura senza
deciderli con regole non dimostrate.

\subparagraph{3. Verifica: istruzioni, dipendenze,
tempi}\label{verifica-istruzioni-dipendenze-tempi}

Codice disponibile (Python 3, libreria standard; ogni script gira in
meno di un minuto): - Nessun codice ancora.

\subparagraph{4. Fonti e contributo}\label{fonti-e-contributo}

Letteratura arXiv (ricerca deterministica
\texttt{tools/cerca\_letteratura.sh}, abstract letti, non usata come
prova): - arXiv:2603.26043v1 --- Finiteness of Disjoint Covering Systems
with Precisely One Repeated Modulus (Yu Hashimoto, 2026); solo abstract
letto. - arXiv:1511.04293v1 --- Searching for Disjoint Covering Systems
with Precisely One Repeated Modulus (Shalosh B. Ekhad, Aviezri S.
Fraenkel, Doron Zeilberger, 2015); solo abstract letto. -
arXiv:2607.24655v1 --- On the problem of large gcd for disjoint residue
classes (Jan Fornal, Yu-Chen Sun, 2026); solo abstract letto. -
arXiv:2608.15873v1 --- Two Questions on \(G\)-harmonic Tuples (Murali
Menon, 2026); solo abstract letto. Come parte 3.

\subparagraph{5. Limiti e parti
irrisolte}\label{limiti-e-parti-irrisolte}

Tutto tranne la riduzione.

\subsubsection{Decisioni degli agenti}
\paragraph{Run \texttt{p4\_c4}}

\subsubsection{Codice su cui poggia il tentativo}
Nessun codice nei tentativi.

\section{arXiv literature consulted}
\begin{thebibliography}{99}
\bibitem{2603.26043v1} Yu Hashimoto. \emph{Finiteness of Disjoint Covering Systems with Precisely One Repeated Modulus}. arXiv:2603.26043v1 (2026). \url{http://arxiv.org/abs/2603.26043v1}
\bibitem{1511.04293v1} Shalosh B. Ekhad, Aviezri S. Fraenkel, Doron Zeilberger. \emph{Searching for Disjoint Covering Systems with Precisely One Repeated Modulus}. arXiv:1511.04293v1 (2015). \url{http://arxiv.org/abs/1511.04293v1}
\bibitem{2607.24655v1} Jan Fornal, Yu-Chen Sun. \emph{On the problem of large gcd for disjoint residue classes}. arXiv:2607.24655v1 (2026). \url{http://arxiv.org/abs/2607.24655v1}
\bibitem{2608.15873v1} Murali Menon. \emph{Two Questions on \$G\$-harmonic Tuples}. arXiv:2608.15873v1 (2026). \url{http://arxiv.org/abs/2608.15873v1}
\end{thebibliography}
\end{document}
